\chapter{Implementation}
\label{ch:implementation}

% ─────────────────────────────────────────────────────────────────────────────
\section{System Architecture}
\label{sec:impl_architecture}
% ─────────────────────────────────────────────────────────────────────────────

GeoMentor is implemented as a monolithic Next.js~14 application deployed to
Vercel's edge network, with Supabase PostgreSQL hosted in the Frankfurt
(eu-central-1) region. The monolithic architecture was a deliberate decision:
the research contribution lies in the adaptive algorithm, not distributed
systems complexity, and a monolith reduces operational overhead and eliminates
inter-service latency. The key architectural constraint is the \textbf{isolation
of the adaptive engine}: the algorithm modules (\texttt{irt.ts} and
\texttt{bkt.ts}) are implemented as pure functions with no direct database
access. All database I/O passes through Next.js API route handlers, which
orchestrate the learner state lifecycle: loading state, invoking engine
functions, persisting updated state, and returning the question. This isolation
makes the engine independently testable and domain-agnostic.

Figure~\ref{fig:component_diagram} shows the four-layer component architecture;
Figure~\ref{fig:sequence_diagram} shows the question lifecycle across both
routes; Figure~\ref{fig:er_diagram} shows the core database schema.

\begin{table}[H]
\centering
\caption{Technology stack with architectural justification.}
\label{tab:tech_stack}
\renewcommand{\arraystretch}{1.25}
\begin{tabularx}{\textwidth}{llX}
\toprule
\textbf{Layer} & \textbf{Technology} & \textbf{Justification} \\
\midrule
Frontend & Next.js 14 (App Router) &
  React Server Components reduce client bundle size; streaming enables
  progressive session summary loading \\
Styling & Tailwind CSS &
  Utility-first CSS; consistent design tokens across all components \\
ORM & Prisma 5 &
  Type-safe database access; migration history provides schema audit trail \\
Database & Supabase PostgreSQL &
  Row-Level Security enforces data isolation at the database layer \\
Auth & Supabase Auth + JWT &
  PKCE flow; bcrypt password hashing; NIST SP 800-63B compliant \\
Deployment & Vercel &
  Preview deployments per branch; automatic HTTPS; edge caching \\
IRT calibration & R + \texttt{mirt} &
  Industry-standard psychometric tooling; offline and audit-logged \\
\bottomrule
\end{tabularx}
\end{table}

The system targets the following non-functional properties: \textbf{scalability}
via Vercel's serverless function infrastructure (each API route invocation is
an independent function execution); \textbf{reliability} via Supabase Frankfurt's
managed PostgreSQL cluster with daily automated backups; \textbf{efficiency}
with EAP grid integration and BKT update executing in-process within the API
route, targeting end-to-end question selection latency under 500\,ms; and
\textbf{maintainability} via TypeScript strict mode, Prisma migration files,
and the adaptive engine's pure-function architecture.

\begin{figure}[H]
\centering
\begin{tikzpicture}[
  box/.style={rectangle, rounded corners=3pt, draw=black!60, fill=white,
              text width=3.3cm, minimum height=0.85cm, align=center, font=\small},
  engine/.style={box, fill=blue!7, draw=blue!45},
  arr/.style={-{Latex[length=2mm]}, thick, black!45, shorten >=2pt, shorten <=2pt},
  lbl/.style={font=\scriptsize, black!55, fill=white, inner sep=1pt}
]
\draw[black!14, dashed] (-5.4, 5.0) -- (5.4, 5.0);
\draw[black!14, dashed] (-5.4, 3.0) -- (5.4, 3.0);
\draw[black!14, dashed] (-5.4, 1.0) -- (5.4, 1.0);

\node[anchor=east, font=\footnotesize\itshape, text=black!40] at (-5.2, 6.0) {Presentation};
\node[anchor=east, font=\footnotesize\itshape, text=black!40] at (-5.2, 4.0) {Application};
\node[anchor=east, font=\footnotesize\itshape, text=black!40] at (-5.2, 2.0) {Adaptive Engine};
\node[anchor=east, font=\footnotesize\itshape, text=black!40] at (-5.2, 0.0) {Persistence};

\node[box] (fe) at (0, 6.0)
  {\textbf{React Frontend}\\{\scriptsize SessionPage $\cdot$ QuestionCard}};
\node[box] (nq) at (-2.8, 4.0)
  {\textbf{GET /next-question}\\{\scriptsize item selection}};
\node[box] (an) at (2.8, 4.0)
  {\textbf{POST /answer}\\{\scriptsize state update}};
\node[engine] (irt) at (-2.8, 2.0)
  {\textbf{IRT Module}\\{\scriptsize \texttt{irt.ts} $\cdot$ pure function}};
\node[engine] (bkt) at (2.8, 2.0)
  {\textbf{BKT Module}\\{\scriptsize \texttt{bkt.ts} $\cdot$ pure function}};
\node[box] (db) at (0, 0.0)
  {\textbf{Supabase PostgreSQL}\\{\scriptsize Prisma ORM $\cdot$ eu-central-1}};

\draw[arr] (fe.south west) -- node[lbl, left]{HTTP GET} (nq.north);
\draw[arr] (fe.south east) -- node[lbl, right]{HTTP POST} (an.north);
\draw[arr] (nq) -- node[lbl, left=2pt]{Fisher info} (irt);
\draw[arr] (an.south west) -- node[lbl, above, sloped]{EAP update} (irt.north east);
\draw[arr] (an) -- node[lbl, right=2pt]{BKT update} (bkt);
\draw[arr] (nq.south east) -- node[lbl, above, sloped]{load variants} (db.north west);
\draw[arr] (an.south west) -- node[lbl, above, sloped]{persist state} (db.north east);
\draw[blue!25, thick, dashed, rounded corners=5pt]
  (-4.5, 1.05) rectangle (4.5, 2.95);
\node[font=\scriptsize\itshape, text=blue!50, anchor=south east] at (4.4, 1.1)
  {zero UI/DB dependencies};
\end{tikzpicture}
\caption{UML 2.5 component diagram. The Adaptive Engine subsystem (IRT and BKT
modules) is implemented as pure functions with no direct UI or database
dependencies.}
\label{fig:component_diagram}
\end{figure}

% ─────────────────────────────────────────────────────────────────────────────
\section{Database Schema}
\label{sec:impl_schema}
% ─────────────────────────────────────────────────────────────────────────────

The Prisma schema defines 20 models. The six central to the adaptive engine
are described here; the full schema is in the repository
(\texttt{prisma/schema.prisma}).

The \textbf{Session} model carries the principal adaptive state: \texttt{theta}
(current EAP estimate, initialised at $-0.780$), \texttt{thetaSd} (posterior
SD, initialised at $0.543$), and \texttt{kcStates} (JSON-serialised
\texttt{Record<kcId, KCState>} carrying BKT state for all 13 KCs). JSON
serialisation eliminates join overhead on the hot adaptive selection path.
The \textbf{Interaction} model records every question--answer event with full
psychometric resolution: \texttt{theta\_before}/\texttt{theta\_after},
\texttt{pLearned\_before}/\texttt{pLearned\_after}, \texttt{hintLevelMax}
(used by the credit factor computation), and \texttt{responseTimeMs}.
The question bank is a two-layer hierarchy: \texttt{QuestionTemplate}
(defines KC, Bloom level, IRT parameters $a$, $b$, $c$, and a stem with
\texttt{\{variable\}} placeholders) and \texttt{QuestionVariant} (surface-specific
substitutions). \texttt{UserVariantSeen} records every (user, variant) pair
presented, enforced as a unique constraint to prevent re-exposure. The
\texttt{AnalyticsEvent} model provides a structured event log including
\texttt{bloom\_escalation}, \texttt{variant\_pool\_exhausted},
\texttt{kc\_mastered}, and \texttt{decay\_detected} events.

\begin{figure}[H]
\centering
\begin{tikzpicture}[
  ent/.style={rectangle, draw=black!60, fill=white, rounded corners=2pt,
              text width=3.2cm, align=center, font=\small, inner sep=4pt},
  hdr/.style={ent, fill=gray!12, font=\small\bfseries, minimum height=0.65cm},
  arr/.style={-{Latex[length=2mm]}, thick, black!50, shorten >=2pt, shorten <=2pt},
  lbl/.style={font=\scriptsize, black!55, fill=white, inner sep=1pt}
]

\node[hdr] (uH) at (0, 8)      {User};
\node[ent, anchor=north] (u) at (uH.south) {
  \texttt{id} (PK)\\
  \texttt{email}\\
  \texttt{xpTotal}\\
  \texttt{kcsMastered}
};

\node[hdr] (sH) at (5, 8)      {Session};
\node[ent, anchor=north] (s) at (sH.south) {
  \texttt{id} (PK)\\
  \texttt{userId} (FK)\\
  \texttt{theta}\\
  \texttt{kcStates} (JSON)\\
  \texttt{condition}
};

\node[hdr] (iH) at (10, 8)     {Interaction};
\node[ent, anchor=north] (i) at (iH.south) {
  \texttt{id} (PK)\\
  \texttt{sessionId} (FK)\\
  \texttt{isCorrect}\\
  \texttt{theta\_pre/post}\\
  \texttt{pLearned\_pre/post}
};

\node[hdr] (tH) at (5, 3.5)    {QuestionTemplate};
\node[ent, anchor=north] (t) at (tH.south) {
  \texttt{id} (PK)\\
  \texttt{kcId}\\
  \texttt{bloomLevel}\\
  \texttt{a}, \texttt{b}, \texttt{c}
};

\node[hdr] (vH) at (10, 3.5)   {QuestionVariant};
\node[ent, anchor=north] (v) at (vH.south) {
  \texttt{id} (PK)\\
  \texttt{templateId} (FK)\\
  \texttt{content}
};

\node[hdr] (uvH) at (5, 0)     {UserVariantSeen};
\node[ent, anchor=north] (uv) at (uvH.south) {
  \texttt{userId} (FK, unique)\\
  \texttt{variantId} (FK, unique)
};

\node[hdr] (susH) at (0, 3.5)  {SUSResponse};
\node[ent, anchor=north] (sus) at (susH.south) {
  \texttt{id} (PK)\\
  \texttt{sessionId} (FK)\\
  \texttt{score}
};

\draw[arr] (uH.east) -- node[lbl, near start, above]{1}
                         node[lbl, near end,   above]{N} (sH.west);
\draw[arr] (sH.east) -- node[lbl, near start, above]{1}
                         node[lbl, near end,   above]{N} (iH.west);
\draw[arr] (tH.east) -- node[lbl, near start, above]{1}
                         node[lbl, near end,   above]{N} (vH.west);
\draw[arr] (vH.north) -- node[lbl, near start, right]{N}
                          node[lbl, near end,   right]{1} (iH.south);
\draw[arr] (u.south) -- ++(0,-0.5) -| (uv.north west)
  node[lbl, near start, right]{1};
\node[lbl] at (3.4, 0.3) {N};
\draw[arr] (v.south) -- node[lbl, near start, right]{1}
                         node[lbl, near end,   right]{N} (uv.north east);
\draw[arr] (sH.west) -- node[lbl, near start, above]{1}
                         node[lbl, near end,   above]{0..1} (susH.east);
\end{tikzpicture}
\caption[ER diagram: six core Prisma schema models.]{Entity-relationship diagram for the six core Prisma schema models.
\texttt{UserVariantSeen} enforces a unique (userId, variantId) constraint,
preventing variant re-exposure across sessions.}
\label{fig:er_diagram}
\end{figure}

% ─────────────────────────────────────────────────────────────────────────────
\section{Adaptive Engine Implementation}
\label{sec:impl_engine}
% ─────────────────────────────────────────────────────────────────────────────

The IRT module (\texttt{src/lib/algorithms/irt.ts}) implements the 3PL model
with the $D = 1.7$ logistic approximation \parencite{lord1980}. The EAP
estimator integrates over a grid of $K = 161$ points spanning $\theta \in
[-4, 4]$ with step $0.05$. The finer grid (compared to the $0.1$-step grid
specified at design time) was adopted after self-pilot testing revealed that
coarser grids produced visible staircase artefacts in the $\hat{\theta}$
trajectory for high-ability learners.

The BKT module (\texttt{src/lib/algorithms/bkt.ts}) implements the standard
Corbett--Anderson update plus the credit factor modification. The implementation
adopts a \textit{soft observation} interpolation: for credit-weighted correct
responses, the posterior is interpolated between the full Bayesian update and
the prior in proportion to $\text{cf}$, which is mathematically equivalent to
a Bayesian update with a degraded likelihood but numerically more stable.
Incorrect responses following hint access receive full-weight Bayesian updates:
providing a hint does not make an incorrect response less informative about
knowledge state.

The BKT mastery threshold is implemented at $p(\text{Learned}) \geq 0.95$,
higher than the $0.80$ threshold cited in the literature review. The binary
nature of BKT's latent state demands a stronger guarantee that the posterior
is not merely elevated by a lucky guess run. The $0.80$ threshold applies
separately to the IRT-based mastery check (\texttt{checkMastery()} in
\texttt{irt.ts}), which averages $P(\text{correct})$ across all KC items at
the current $\hat{\theta}$ using the continuous IRT scale.

% ─────────────────────────────────────────────────────────────────────────────
\section{API Route Design}
\label{sec:impl_api}
% ─────────────────────────────────────────────────────────────────────────────

The adaptive session is managed by four API routes. \textbf{GET
/api/sessions/[id]/next-question}: loads session state; builds the eligible
variant pool (joining \texttt{QuestionTemplate} and \texttt{QuestionVariant}
minus \texttt{UserVariantSeen}); applies mastery gate, Bloom gate, deduplication,
and decay boost; selects $i^* = \arg\max I_i(\hat{\theta})$; renders variant
by interpolating placeholders and shuffling answer options (Fisher-Yates);
returns serialised question object. \textbf{POST /api/sessions/[id]/answer}:
validates the submitted answer; computes credit factor from \texttt{hintLevelMax};
runs EAP update and BKT update; persists updated session state atomically;
writes \texttt{Interaction} record; checks mastery and Bloom escalation
thresholds; emits analytics events as appropriate.
\textbf{GET /api/sessions/[id]/summary}: generates the five-tab dashboard payload.
\textbf{GET /api/sessions/[id]/trajectory}: returns the longitudinal
$\hat{\theta}$ sequence for the ability progression chart. Figure~\ref{fig:sequence_diagram}
illustrates the complete question lifecycle.

\begin{figure}[H]
\centering
\begin{tikzpicture}[
  actor/.style={rectangle, rounded corners=2pt, draw=black!60, fill=gray!8,
                text width=1.85cm, minimum height=0.65cm, align=center,
                font=\footnotesize\bfseries},
  msg/.style={-{Latex[length=2mm]}, thick, black!55, shorten >=2pt},
  ret/.style={-{Latex[length=2mm]}, thick, black!35, dashed, shorten >=2pt},
  lbl/.style={font=\scriptsize, black!60, fill=white, inner sep=1pt}
]
\def\xL{0}
\def\xF{2.2}
\def\xA{4.8}
\def\xE{7.6}
\def\xD{10.8}

\node[actor] at (\xL, 9.5) {Learner};
\node[actor] at (\xF, 9.5) {Frontend};
\node[actor] at (\xA, 9.5) {API Routes};
\node[actor] at (\xE, 9.5) {IRT + BKT\\Engine};
\node[actor] at (\xD, 9.5) {Supabase\\PostgreSQL};

\draw[black!18, dashed] (\xL, 9.1) -- (\xL, 0.3);
\draw[black!18, dashed] (\xF, 9.1) -- (\xF, 0.3);
\draw[black!18, dashed] (\xA, 9.1) -- (\xA, 0.3);
\draw[black!18, dashed] (\xE, 9.1) -- (\xE, 0.3);
\draw[black!18, dashed] (\xD, 9.1) -- (\xD, 0.3);

\node[anchor=west, font=\scriptsize\itshape, text=black!40] at (-0.5, 8.85)
  {POST /answer};

\draw[msg] (\xL, 8.6) -- (\xF, 8.6)
  node[lbl, midway, above]{submit response};
\draw[msg] (\xF, 8.1) -- (\xA, 8.1)
  node[lbl, midway, above]{POST /answer};
\draw[msg] (\xA, 7.5) -- (\xD, 7.5)
  node[lbl, midway, above]{load state};
\draw[ret] (\xD, 7.1) -- (\xA, 7.1)
  node[lbl, midway, above]{theta, kcStates};
\draw[msg] (\xA, 6.5) -- (\xE, 6.5)
  node[lbl, midway, above]{EAP + BKT update};
\draw[ret] (\xE, 6.1) -- (\xA, 6.1)
  node[lbl, midway, above]{theta\_new, pLearned\_new};
\draw[msg] (\xA, 5.5) -- (\xD, 5.5)
  node[lbl, midway, above]{persist state};
\draw[ret] (\xD, 5.1) -- (\xA, 5.1)
  node[lbl, midway, above]{ack};
\draw[ret] (\xA, 4.5) -- (\xF, 4.5)
  node[lbl, midway, above]{feedback payload};
\draw[ret] (\xF, 4.1) -- (\xL, 4.1)
  node[lbl, midway, above]{render feedback};

\draw[black!15] (-0.5, 3.75) -- (11.5, 3.75);
\node[anchor=west, font=\scriptsize\itshape, text=black!40] at (-0.5, 3.6)
  {GET /next-question};

\draw[msg] (\xF, 3.4) -- (\xA, 3.4)
  node[lbl, midway, above]{GET /next-question};
\draw[msg] (\xA, 2.9) -- (\xD, 2.9)
  node[lbl, midway, above]{load unseen variants};
\draw[ret] (\xD, 2.5) -- (\xA, 2.5)
  node[lbl, midway, above]{eligible pool};
\draw[msg] (\xA, 2.0) -- (\xE, 2.0)
  node[lbl, midway, above]{Fisher select};
\draw[ret] (\xE, 1.6) -- (\xA, 1.6)
  node[lbl, midway, above]{next variant};
\draw[ret] (\xA, 1.1) -- (\xF, 1.1)
  node[lbl, midway, above]{question payload};
\draw[ret] (\xF, 0.6) -- (\xL, 0.6)
  node[lbl, midway, above]{render question};
\end{tikzpicture}
\caption{Sequence diagram of the adaptive question lifecycle across both API
routes. The IRT and BKT engine functions are pure computations invoked
synchronously; all database I/O is mediated by the route layer. Target
round-trip latency: $<500$\,ms.}
\label{fig:sequence_diagram}
\end{figure}

% ─────────────────────────────────────────────────────────────────────────────
\section{Security Implementation}
\label{sec:impl_security}
% ─────────────────────────────────────────────────────────────────────────────

Security requirements were specified against NIST SP 800-63B \parencite{nist2017}:

\begin{description}
  \item[Password storage.] bcrypt at cost factor 12; plaintext never persisted.
  \item[Authentication.] JWT with Supabase Auth PKCE flow; tokens short-lived (1 hour); payloads contain user ID and role only.
  \item[Rate limiting.] Auth endpoints rate-limited to 10 requests/minute/IP, mitigating credential stuffing.
  \item[Row-Level Security.] Supabase RLS enforces per-user data isolation on all 33 tables; verified by cross-user access testing.
  \item[Input validation.] All API route inputs validated server-side against typed schemas; client-side validation supplementary.
  \item[HTTPS.] All traffic served over HTTPS via Vercel edge certificates; HTTP redirected at CDN layer.
\end{description}

% ─────────────────────────────────────────────────────────────────────────────
\section{Frontend Components}
\label{sec:impl_frontend}
% ─────────────────────────────────────────────────────────────────────────────

The learner session follows a linear state machine: \texttt{pre-test}
$\to$ \texttt{instruction} $\to$ \texttt{question} $\to$ \texttt{feedback}
$\to$ \texttt{summary}. The question interface renders the stem and four
shuffled answer options (Fisher-Yates to prevent position bias), a progress
bar for session and KC mastery, a three-level hint panel with progressive
disclosure, a confidence slider, and a pause button that persists session
state to the database for future resumption.

\texttt{FirstVisitTour.tsx} implements a five-step onboarding overlay explaining
the adaptive engine, the BKT mastery bar, the hint system, the XP level
system, and the session summary. Tour completion is persisted to
\texttt{localStorage}. This implements Mayer's (2009) signalling principle:
the adaptive system's reasoning is made legible before the learner encounters
the first question.

\texttt{SessionSummaryDashboard.tsx} presents a five-tab post-session analysis
interface: Overview (statistics and performance narrative), Trajectory
($\hat{\theta}$ chart with hover tooltips), Skills (per-KC mastery bar chart),
Misconceptions (distractor frequency table), and Share (PNG performance card
export). The leaderboard ranks opt-in learners by total XP; consent is
double-checked server-side on every read.

Figure~\ref{fig:ss_dashboard} and Figure~\ref{fig:ss_hint} show the live
production interface for two key interaction states.

\begin{figure}[H]
\centering
\includegraphics[width=0.82\textwidth]{ss_dashboard}
\caption{GeoMentor dashboard (hi-fi screenshot): adaptive badge, XP display,
KC mastery bar, and session entry point. The ``Adaptive'' badge communicates
the personalisation mode to the learner.}
\label{fig:ss_dashboard}
\end{figure}

\begin{figure}[H]
\centering
\includegraphics[width=0.82\textwidth]{ss_hint}
\caption{Three-level hint panel (hi-fi screenshot): conceptual hint displayed
at Level~1, with progressive disclosure controls for Level~2 (procedural
strategy) and Level~3 (worked example). The BKT credit factor is reduced for
any correct answer following hint access.}
\label{fig:ss_hint}
\end{figure}

% ─────────────────────────────────────────────────────────────────────────────
\section{Challenges Encountered}
\label{sec:impl_challenges}
% ─────────────────────────────────────────────────────────────────────────────

\textbf{EAP grid resolution.} The specified 81-point grid produced visible
staircase discontinuities in $\hat{\theta}$ trajectories for high-ability
learners. The grid was refined to 161 points (step $0.05$), resolving the
artefact at acceptable computational cost.

\textbf{Review Material bug.} The instruction mode interstitial triggers
an unintended \texttt{setAppState('start')} call in certain render paths,
resetting the session to the start screen. This affects only the post-mastery
instructional interstitial, not the adaptive engine's core question selection
or BKT update logic. The bug is disclosed here in full; a fix was deprioritised
in Sprint~9 in favour of evaluation data collection.

\textbf{K-means cold-start: evaluated and superseded.} The interim report
described K-means clustering for cold-start initialisation. Following
implementation of pre-test seeding, K-means was evaluated and superseded:
pre-test seeding provides direct per-KC ability evidence rather than a
demographic proxy, is more interpretable, and avoids conditioning initial
learning experience on cluster assignment.

% ─────────────────────────────────────────────────────────────────────────────
\section{Requirements Traceability}
\label{sec:impl_traceability}
% ─────────────────────────────────────────────────────────────────────────────

\begin{table}[H]
\centering
\caption{Requirements traceability matrix: functional requirements to
implementation artefacts.}
\label{tab:traceability}
\renewcommand{\arraystretch}{1.2}
\begin{tabularx}{\textwidth}{lXlX}
\toprule
\textbf{Req.} & \textbf{Requirement} & \textbf{Status} & \textbf{Artefact} \\
\midrule
FR01 & Adaptive question selection using IRT 3PL & Done &
  \texttt{irt.ts:selectNextQuestion} \\
FR02 & Per-KC knowledge state tracking (BKT) & Done &
  \texttt{bkt.ts:updateBKT} \\
FR03 & Bloom-gated question escalation & Done &
  \texttt{next-question/route.ts} \\
FR04 & Pre-test ability seeding & Done &
  \texttt{irt.ts:estimateTheta} \\
FR05 & Cross-session variant deduplication & Done &
  \texttt{UserVariantSeen} model \\
FR06 & Graduated hint system with credit penalty & Done &
  \texttt{hints/route.ts}; \texttt{bkt.ts:updateBKT} \\
FR07 & BKT decay detection + successive relearning & Done &
  \texttt{next-question/route.ts} (decay boost) \\
FR08 & Session persistence and resume & Done &
  \texttt{Session.kcStates} JSON field \\
FR09 & SUS usability instrument & Done &
  Post-session questionnaire \\
FR10 & Leaderboard with consent opt-in & Done &
  \texttt{leaderboard/route.ts}; \texttt{User.leaderboardOptIn} \\
FR11 & NIST 800-63B compliant authentication & Done &
  Supabase Auth + bcrypt + rate limiting \\
FR12 & Analytics event log & Done &
  \texttt{AnalyticsEvent} model \\
FR13 & 7-day retention assessment & Partial &
  Infrastructure implemented; data collection limited by timeline \\
FR14 & Multi-domain KC support & Deferred &
  Architecture supports; content not authored \\
\bottomrule
\end{tabularx}
\end{table}

% ─────────────────────────────────────────────────────────────────────────────
\section{Deployment, Testing, and Quality}
\label{sec:impl_ci}
% ─────────────────────────────────────────────────────────────────────────────

The build pipeline executes on every push to main: \texttt{prisma generate}
(type safety); \texttt{tsc --noEmit} (TypeScript strict-mode check, fails
on any type error); \texttt{next build} (production build); and Vercel
deployment to preview or production URL.

Three testing layers were implemented. \textbf{Unit tests}: 79 Jest tests
cover the IRT and BKT algorithm modules, validating the 3PL probability
function at boundary values, EAP convergence on synthetic sequences, BKT
posterior monotonicity, credit-factor interpolation, and Bloom gate enforcement.
The pure-function architecture makes these directly testable without mocking.
\textbf{Integration testing}: a structured developer self-pilot on 2026-02-26
verified the end-to-end session flow and provided the empirical prior parameters
($\hat{\theta}_0 = -0.780$, $\sigma_0 = 0.543$). \textbf{Regression testing}:
TypeScript strict mode and Next.js production build execute automatically on
every push; feature branches are deployed to isolated Vercel preview URLs
before merge. The primary quality gap is the absence of automated end-to-end
browser tests; UI regression relied on manual verification against preview
deployments, disclosed per the CSSE quality assurance principle.
